\section{Experiments}

The experiments test whether proprioceptive self-grounding improves spatial
reliability beyond ordinary action supervision and state concatenation. The
evaluation links three evidence types: task performance, action stability under
camera perturbations, and visual grounding of the acting robot.

\subsection{Experimental Setup}

We evaluate a VLA backbone fine-tuned on language-conditioned manipulation
demonstrations in simulation and on a real robot. Each episode contains visual
observations \(o_t\), proprioceptive state \(s_t\), language instruction \(x\),
and action supervision \(a_t\). For robustness evaluation, we keep the physical
robot-scene state, instruction, and proprioceptive state fixed, then apply
lateral view shifts of controlled magnitude. This isolates camera-induced
changes from changes in the underlying manipulation state.

We compare four policies under the same demonstrations, action labels, and
view-shift protocol: Base VLA, action-supervised fine-tuning, VLA with
proprioception concatenation, and \textsc{ProprioVLA}. We report task success
rate, action prediction error, Trajectory Shift Ratio, Shift Direction
Agreement, and a grounding score. The grounding score is computed as either
attention-heatmap alignment with \(H_t\) or attention mass on robot regions.
Lower trajectory shift and direction agreement indicate a more stable action
frame; higher grounding score indicates stronger localization of the acting
body. Each setup reports the task suite, demonstration count, train/test split,
camera shifts, action parameterization, proprioceptive normalization,
optimizer, and loss weights \(\lambda_g\), \(\lambda_r\), and \(\beta\).

\subsection{Simulation Experiments}

Simulation provides controlled access to camera shifts, object layouts, and
ground-truth robot geometry. We evaluate normal, mild-shift, medium-shift, and
large-shift settings. For each setting, we report task success, action
prediction error, Trajectory Shift Ratio, Shift Direction Agreement, and
grounding score. Robustness curves plot these metrics against shift magnitude,
making the co-movement between predicted actions and image shifts directly
visible.

We also evaluate held-out object layouts and camera positions. These splits
separate memorization of a fixed training camera from grounding that transfers
across views. Qualitative trajectory overlays compare the original prediction
with shifted-view predictions after mapping them to the same robot-centric
coordinate frame.

\subsection{Real-World Experiments}

The real-robot evaluation tests whether the same failure mode appears under
sensor noise, calibration error, and real scene variation. We evaluate several
manipulation tasks under the original view and shifted views, reporting task
success, success drop under perturbation, action prediction error on logged
demonstrations, and grounding score.

For qualitative analysis, we visualize baseline attention, the mined
state-aligned heatmap, \textsc{ProprioVLA} attention, and trajectory overlays
for the same scene. The comparison checks whether attention remains dominated
by task objects and background context or shifts toward the arm, gripper, and
end-effector while the trajectory remains stable.

\subsection{Ablations}

The ablation study isolates the mechanisms that distinguish
\textsc{ProprioVLA} from simple state conditioning. We evaluate four variants:
(1) removing the visual self-grounding loss \(\mathcal{L}_{\mathrm{ground}}\);
(2) removing the reconstruction loss \(\mathcal{L}_{\mathrm{rec}}\); (3) using
raw optical flow instead of state-aligned residual motion; and (4) replacing
the proprioceptive visual query with state concatenation. Each ablation is
evaluated with the same success, action, trajectory-shift, and grounding
metrics.

These ablations test the central mechanism: the policy must use proprioception
to ground visual attention and anchor the attended feature, rather than merely
receive proprioceptive values as extra inputs. The comparison separates three
alternatives: gains from additional state, gains from motion saliency alone,
and gains from explicit proprioceptive self-grounding.
